% Môn: Vật lý
% Lớp: 11
% Chương: Chương I. Dao động
% Bài: L11C1 Bài 2. Mô tả dao động điều hòa · Xác định các đại lượng đặc trưng từ phương trình li độ
% ===== Câu 44 =====
% ID: L11C1B2-131-DS
% Mức: TH
\dangbt{Xác định các đại lượng đặc trưng từ phương trình li độ}
\begin{ex}
% ID: L11C1B2-131-DS
% Mức: TH
Một vật dao động điều hòa dọc theo trục Ox có phương trình li độ $x = 4\cos\left(5t - \dfrac{\pi}{2}\right)$ (cm), trong đó $t$ tính bằng giây. Xét tính chính xác của các mệnh đề sau:
\choiceTF
{\True Tần số góc của dao động là $5$ rad/s.}
{\True Pha của dao động ở thời điểm $t$ bất kì là $\left(5t - \dfrac{\pi}{2}\right)$ rad.}
{Tại thời điểm ban đầu $t = 0$, vật đi qua vị trí cân bằng theo chiều âm.}
{\True Chu kì dao động của vật là $0,4\pi$ s.}
\loigiai{
\begin{itemchoice}
\item (Đúng) Tần số góc của dao động là $5$ rad/s. (Vì): Mệnh đề thứ nhất đúng vì tần số góc của dao động là $\omega = 5$ rad/s
\item (Đúng) Pha của dao động ở thời điểm $t$ bất kì là $\left(5t - \dfrac{\pi}{2}\right)$ rad. (Vì): Mệnh đề thứ hai đúng vì pha của dao động ở thời điểm $t$ là toàn bộ biểu thức góc trong hàm cosin: $\Phi(t) = 5t - \dfrac{\pi}{2}$ rad
\item (Sai) Tại thời điểm ban đầu $t = 0$, vật đi qua vị trí cân bằng theo chiều âm. (Vì): Mệnh đề thứ ba sai vì tại $t = 0$, pha ban đầu là $\varphi = -\dfrac{\pi}{2}$ rad, tương ứng với vật đi qua vị trí cân bằng theo chiều dương ($v \gt  0$
\item (Đúng) Chu kì dao động của vật là $0,4\pi$ s. (Vì): Mệnh đề thứ tư đúng vì chu kì dao động được tính theo công thức: $T = \dfrac{2\pi}{\omega} = \dfrac{2\pi}{5} = 0,4\pi$ s
\end{itemchoice}
}
\end{ex}

% ===== Câu 46 =====
% ID: L11C1B2-132-DS
% Mức: TH
\begin{ex}
% ID: L11C1B2-132-DS
% Mức: TH
Một chất điểm dao động điều hòa có phương trình li độ biểu diễn qua hàm sin là $x = 6\sin\left(2\pi t + \dfrac{\pi}{3}\right)$ (cm). Đánh giá các nhận định sau đây:
\choiceTF
{\True Biên độ dao động của chất điểm bằng $6$ cm.}
{Pha ban đầu của dao động khi chuyển đổi về dạng hàm cosin chuẩn dương là $\dfrac{\pi}{3}$ rad.}
{\True Chu kì dao động của chất điểm là $1$ s.}
{Tần số dao động của chất điểm là $2\pi$ Hz.}
\loigiai{
\begin{itemchoice}
\item (Đúng) Biên độ dao động của chất điểm bằng $6$ cm. (Vì): Mệnh đề thứ nhất đúng vì biên độ dao động của chất điểm là $A = 6$ cm
\item (Sai) Pha ban đầu của dao động khi chuyển đổi về dạng hàm cosin chuẩn dương là $\dfrac{\pi}{3}$ rad. (Vì): Mệnh đề thứ hai sai vì để đưa phương trình về dạng hàm cosin chuẩn, ta dùng công thức $\sin\alpha = \cos\left(\alpha - \dfrac{\pi}{2}\right)$, khi đó: $x = 6\cos\left(2\pi t + \dfrac{\pi}{3} - \dfrac{\pi}{2}\right) = 6\cos\left(2\pi t - \dfrac{\pi}{6}\right)$
\item (Đúng) Chu kì dao động của chất điểm là $1$ s. (Vì): Mệnh đề thứ ba đúng vì chu kì dao động của chất điểm là $T = \dfrac{2\pi}{\omega} = \dfrac{2\pi}{2\pi} = 1$ s
\item (Sai) Tần số dao động của chất điểm là $2\pi$ Hz.
\end{itemchoice}
}
\end{ex}

% ===== Câu 50 =====
% ID: L11C1B2-133-DS
% Mức: TH
\begin{ex}
% ID: L11C1B2-133-DS
% Mức: TH
Một chất điểm dao động điều hòa với phương trình li độ $x = 10\cos\left(4\pi t - \dfrac{\pi}{3}\right)$ (cm), trong đó thời gian $t$ tính bằng giây. Xét tính đúng/sai của các phát biểu sau:
\choiceTF
{\True Biên độ dao động của chất điểm là $10$ cm.}
{Chu kì dao động của chất điểm là $2$ s.}
{\True Tần số dao động của chất điểm là $2$ Hz.}
{Pha ban đầu của dao động là $\dfrac{\pi}{3}$ rad.}
\loigiai{
\begin{itemchoice}
\item (Đúng) Biên độ dao động của chất điểm là $10$ cm. (Vì): Mệnh đề thứ nhất đúng vì đối chiếu với phương trình tổng quát $x = A\cos(\omega t + \varphi)$, ta có hệ số biên độ $A = 10$ cm
\item (Sai) Chu kì dao động của chất điểm là $2$ s. (Vì): Mệnh đề thứ hai sai vì chu kì dao động của vật được tính bằng công thức $T = \dfrac{2\pi}{\omega} = \dfrac{2\pi}{4\pi} = 0,5$ s
\item (Đúng) Tần số dao động của chất điểm là $2$ Hz. (Vì): Mệnh đề thứ ba đúng vì tần số dao động của vật là $f = \dfrac{1}{T} = \dfrac{1}{0,5} = 2$ Hz
\item (Sai) Pha ban đầu của dao động là $\dfrac{\pi}{3}$ rad. (Vì): Mệnh đề thứ tư sai vì pha ban đầu của dao động là hệ số tự do trong hàm cosin, tức là $\varphi = -\dfrac{\pi}{3}$ rad
\end{itemchoice}
}
\end{ex}

% ===== Câu 52 =====
% ID: L11C1B2-134-TLN
% Mức: VD
\begin{ex}
% ID: L11C1B2-134-TLN
% Mức: VD
[SGK CTST]
	Một vật dao động điều hòa với biên độ $A=6~\mathrm{cm}$ và tần số góc $\omega=5~\mathrm{rad/s}$. Tính tốc độ cực đại của vật (làm tròn đến số nguyên gần nhất, đơn vị cm/s).
\shortans{29/01/1900}
\loigiai{
$v_{\max}=A\omega=6\cdot 5=30\ \mathrm{cm/s}$.
}
\end{ex}

% ===== Câu 54 =====
% ID: L11C1B2-135-DS
% Mức: TH
\begin{ex}
% ID: L11C1B2-135-DS
% Mức: TH
Cho các phát biểu sau về dao động, phát biểu nào đúng, phát biểu nào sai?
\choiceTF
{\True Dao động tuần hoàn đơn giản nhất là dao động điều hòa}
{\True Dao động điều hòa là dao động trong đó li độ của vật là một hàm hay sin hay hàm tan theo thời gian}
{\True Phương trình được mô tả dưới dạng $x=A\cos (\omega t+\varphi)$ được gọi là phương trình dao động điều hòa}
{\True Đại lượng $\omega t + \varphi$ được gọi là pha ban đầu của dao động}
\loigiai{
\begin{itemchoice}
\item (Đúng) Dao động tuần hoàn đơn giản nhất là dao động điều hòa. (Vì): ao động tuần hoàn đơn giản nhất là dao động điều hòa
\item (Đúng) Dao động điều hòa là dao động trong đó li độ của vật là một hàm hay sin hay hàm tan theo thời gian. (Vì): ao động điều hòa có li độ là hàm sin hoặc cos, không phải là hàm tan
\item (Đúng) Phương trình được mô tả dưới dạng $x=A\cos (\omega t+\varphi)$ được gọi là phương trình dao động điều hòa. (Vì): Phương trình $x=A\cos (\omega t + \varphi)$ là phương trình của dao động điều hòa
\item (Đúng) Đại lượng $\omega t + \varphi$ được gọi là pha ban đầu của dao động. (Vì): ại lượng $(\omega t + \varphi)$ là pha dao động tại thời điểm $t$, còn pha ban đầu là $\varphi$
\end{itemchoice}
}
\end{ex}

% ===== Câu 55 =====
% ID: L11C1B2-136-TLN
% Mức: VD
\begin{ex}
% ID: L11C1B2-136-TLN
% Mức: VD
Một vật dao động điều hòa với biên độ $5\,\text{cm}$ và tần số góc $8\,\text{rad/s}$. Tính vận tốc cực đại ($v_{max}$) (cm/s).
\shortans{08/02/1900}
\loigiai{
$v_{max} = \omega A = 8 \cdot 5 = 40.$
}
\end{ex}

% ===== Câu 57 =====
% ID: L11C1B2-137-DS
% Mức: TH
\begin{ex}
% ID: L11C1B2-137-DS
% Mức: TH
Xét phương trình li độ tổng quát của một vật dao động điều hòa $x = A\cos(\omega t + \varphi)$. Hãy nhận xét các phát biểu dưới đây về mặt lí thuyết:
\choiceTF
{Biên độ $A$ có thể nhận giá trị âm hoặc dương tùy thuộc vào cách chọn gốc tọa độ và gốc thời gian.}
{\True Tần số góc $\omega$ luôn là một đại lượng có giá trị dương.}
{\True Pha ban đầu $\varphi$ là đại lượng dùng để xác định trạng thái ban đầu của vật ở thời điểm $t = 0$.}
{Đơn vị đo của pha dao động $(\omega t + \varphi)$ trong hệ SI là héc (Hz).}
\loigiai{
\begin{itemchoice}
\item (Sai) Biên độ $A$ có thể nhận giá trị âm hoặc dương tùy thuộc vào cách chọn gốc tọa độ và gốc thời gian. (Vì): Mệnh đề thứ nhất sai vì biên độ $A$ là độ lệch cực đại của vật khỏi vị trí cân bằng nên luôn mang giá trị dương ($A \gt  0$
\item (Đúng) Tần số góc $\omega$ luôn là một đại lượng có giá trị dương. (Vì): Mệnh đề thứ hai đúng vì tần số góc $\omega$ đặc trưng cho tốc độ biến thiên của pha dao động và luôn xác định là số dương
\item (Đúng) Pha ban đầu $\varphi$ là đại lượng dùng để xác định trạng thái ban đầu của vật ở thời điểm $t = 0$. (Vì): Mệnh đề thứ ba đúng vì tại thời điểm $t = 0$, li độ và chiều chuyển động của vật được xác định hoàn toàn qua giá trị của pha ban đầu $\varphi$
\item (Sai) Đơn vị đo của pha dao động $(\omega t + \varphi)$ trong hệ SI là héc (Hz). (Vì): Mệnh đề thứ tư sai vì đơn vị của pha dao động trong hệ SI là radian (rad), còn héc (Hz) là đơn vị của tần số $f$
\end{itemchoice}
}
\end{ex}

% ===== Câu 59 =====
% ID: L11C1B2-138-TLN
% Mức: VD
\begin{ex}
% ID: L11C1B2-138-TLN
% Mức: VD
[SGK CTST]
	Một vật dao động điều hòa với biên độ $A=3~\mathrm{cm}$ và tần số góc $\omega=4~\mathrm{rad/s}$. Tính gia tốc cực đại của vật (đơn vị $~\mathrm{cm/s^2}$).
\shortans{16/02/1900}
\loigiai{
$a_{\max}=A\omega^2=3\cdot 4^2=3\cdot 16=48\ \mathrm{cm/s^2}$.
}
\end{ex}

% ===== Câu 64 =====
% ID: L11C1B2-139-DS
% Mức: TH
\begin{ex}
% ID: L11C1B2-139-DS
% Mức: TH
Một vật nhỏ dao động điều hòa theo phương trình có dạng chưa chuẩn $x = -5\cos\left(10\pi t + \dfrac{\pi}{6}\right)$ (cm). Khảo sát các thông số đặc trưng của dao động này:
\choiceTF
{Biên độ dao động của vật là $-5$ cm.}
{\True Phương trình dao động của vật khi viết dưới dạng chuẩn biên độ dương là $x = 5\cos\left(10\pi t - \dfrac{5\pi}{6}\right)$ (cm).}
{\True Tần số góc của dao động bằng $10\pi$ rad/s.}
{Pha ban đầu của dao động sau khi chuẩn hóa là $\dfrac{\pi}{6}$ rad.}
\loigiai{
\begin{itemchoice}
\item (Sai) Biên độ dao động của vật là $-5$ cm. (Vì): Mệnh đề thứ nhất sai vì biên độ dao động $A$ luôn là một đại lượng dương, do đó $A = 5$ cm
\item (Đúng) Phương trình dao động của vật khi viết dưới dạng chuẩn biên độ dương là $x = 5\cos\left(10\pi t - \dfrac{5\pi}{6}\right)$ (cm). (Vì): Mệnh đề thứ hai đúng vì ta áp dụng công thức lượng giác $-\cos\alpha = \cos(\alpha - \pi)$ để khử dấu trừ trước hàm số: $x = -5\cos\left(10\pi t + \dfrac{\pi}{6}\right) = 5\cos\left(10\pi t + \dfrac{\pi}{6} - \pi\right) = 5\cos\left(10\pi t - \dfrac{5\pi}{6}\right)$
\item (Đúng) Tần số góc của dao động bằng $10\pi$ rad/s. (Vì): Mệnh đề thứ ba đúng vì tần số góc của dao động là hệ số của thời gian $t$, tức là $\omega = 10\pi$ rad/s
\item (Sai) Pha ban đầu của dao động sau khi chuẩn hóa là $\dfrac{\pi}{6}$ rad.
\end{itemchoice}
}
\end{ex}

% ===== Câu 88 =====
% ID: L11C1B2-140-TN
% Mức: NB
\begin{ex}
% ID: L11C1B2-140-TN
% Mức: NB
Một vật dao động điều hòa với phương trình $x = 6\sin(5t + \frac{\pi}{4})$ cm. Pha ban đầu của dao động là bao nhiêu?
\choice
{$\frac{\pi}{2}$ rad}
{$\frac{\pi}{3}$ rad}
{$\frac{\pi}{6}$ rad}
{\True $\frac{\pi}{4}$ rad}
\loigiai{
Phương trình dao động điều hòa có dạng $x = A\sin(\omega t + \varphi)$. Pha ban đầu là $\varphi = \frac{\pi}{4}$ rad.
}
\end{ex}

% ===== Câu 90 =====
% ID: L11C1B2-141-TN
% Mức: NB
\begin{ex}
% ID: L11C1B2-141-TN
% Mức: NB
Một vật dao động điều hòa với phương trình $x = 3\sin(4\pi t)$ cm. Tần số góc của dao động là bao nhiêu?
\choice
{$2\pi$ rad/s}
{\True $4\pi$ rad/s}
{$8\pi$ rad/s}
{$\pi$ rad/s}
\loigiai{
Phương trình dao động điều hòa có dạng $x = A\sin(\omega t)$. Tần số góc là $\omega = 4\pi$ rad/s.
}
\end{ex}

% ===== Câu 94 =====
% ID: L11C1B2-142-TN
% Mức: NB
\begin{ex}
% ID: L11C1B2-142-TN
% Mức: NB
Trong biểu thức dao động điều hòa $x = A\cos(\omega t + \varphi)$, đại lượng $\omega$ được gọi là
\choice
{Tần số dao động}
{\True Tần số góc}
{Chu kì dao động}
{Biên độ dao động}
\loigiai{
Tần số góc $\omega$ liên hệ với chu kì bởi $\omega = \dfrac{2\pi}{T}$, đơn vị là rad/s.
}
\end{ex}

% ===== Câu 96 =====
% ID: L11C1B2-143-TN
% Mức: NB
\begin{ex}
% ID: L11C1B2-143-TN
% Mức: NB
Biết pha ban đầu của một vật dao động điều hòa, ta xác định được
\choice
{cách kích thích dao động}
{chu kỳ và trạng thái dao động}
{\True chiều chuyển động của vật lúc ban đầu}
{quỹ đạo dao động}
\loigiai{
Pha ban đầu cho biết trạng thái dao động ban đầu: vị trí và chiều chuyển động của vật tại $t = 0$.
}
\end{ex}

% ===== Câu 100 =====
% ID: L11C1B2-144-TN
% Mức: NB
\begin{ex}
% ID: L11C1B2-144-TN
% Mức: NB
$\immini$ {
 Hình vẽ mô tả một vật bắt đầu dao động điều hòa tại vị trí lệch pha so với gốc thời gian. Đại lượng $\varphi)$ trong phương trình $x = A\cos(\omega t + \varphi)$ gọi là gì?
\choice
{Biên độ dao động}
{Chu kì dao động}
{\True Pha ban đầu}
{Tần số dao động}
\loigiai{
Pha ban đầu $\varphi$ xác định trạng thái ban đầu của vật tại $t = 0$.
}
\end{ex}

% ===== Câu 104 =====
% ID: L11C1B2-145-TN
% Mức: NB
\begin{ex}
% ID: L11C1B2-145-TN
% Mức: NB
Một chất điểm dao động điều hòa với phương trình $x=A\cos (\omega t+\varphi)$, trong đó $\omega$ có giá trị dương. Đại lượng $\omega$ gọi là:
\choice
{Biên độ dao động}
{Chu kì của dao động}
{\True Tần số góc của dao động}
{Pha ban đầu của dao động}
\loigiai{
Đại lượng $\omega$ là tần số góc của dao động, đơn vị là $\mathrm{rad} / s$
}
\end{ex}

% ===== Câu 108 =====
% ID: L11C1B2-146-TN
% Mức: NB
\begin{ex}
% ID: L11C1B2-146-TN
% Mức: NB
Độ lệch cực đại so với vị trí cân bằng gọi là
\choice
{\True Biên độ}
{Tần số}
{Li độ}
{Pha ban đầu}
\loigiai{
Biên độ là độ lệch lớn nhất của vật dao động so với vị trí cân bằng
}
\end{ex}

% ===== Câu 110 =====
% ID: L11C1B2-147-TN
% Mức: NB
\begin{ex}
% ID: L11C1B2-147-TN
% Mức: NB
Tần số góc có đơn vị là
\choice
{Hz}
{cm}
{rad}
{\True rad / s}
\loigiai{
Đơn vị của tần số góc là radian trên giây ($\mathrm{rad} / s$)
}
\end{ex}
